\documentclass[11pt,a4paper]{article}
\usepackage[T1]{fontenc}
\usepackage{lmodern,amsmath,amssymb,amsthm,booktabs,microtype}
\usepackage[margin=26mm]{geometry}
\usepackage[colorlinks=true]{hyperref}
\newtheorem{proposition}{Proposition}
\newtheorem{corollary}[proposition]{Corollary}
\newcommand{\Q}{\mathbb Q}\newcommand{\Z}{\mathbb Z}
\DeclareMathOperator{\Tr}{Tr}
\title{Fixed multiples on the small trace cubic and actual detector contacts}
\author{}\date{17 November 2026}
\begin{document}\maketitle
\begin{abstract}
We compute the nodal specializations of all fixed multiples of the
reference point on the small trace cubic. A finite menu of these
sections and their two-torsion translates cannot realize the actual
detector traces in the large-slope preparation. It also fails at every
fresh contact outside an explicitly determined finite prime set.
The statements do not exclude individual rational witnesses or
indices varying with the source tuple.
\end{abstract}

\section{The nodal specialization}
Consider
\[
 E_w:\ y^2=x\bigl(x^2-(w+4)x+4(w+1)\bigr),\qquad
 P=(4,4),\quad T=(0,0),\quad V=x-2.
\]
For $m\ge1$ define the positive integers $a_m,b_m$ by
\[
 (3+2\sqrt2)^m=a_m+b_m\sqrt2.
\]
Thus $a_m^2-2b_m^2=1$, and
\[
 a_1=3,\quad b_1=2,\qquad
 a_{m+1}=3a_m+4b_m,\quad b_{m+1}=2a_m+3b_m.
\]
Write $V_m(w)=x([m]P)-2$ and
$V_{m,T}(w)=x([m]P+T)-2$.

\begin{proposition}\label{prop:limits}
For every nonzero fixed integer $m$, both rational functions are
regular at $w=0$, and
\[
 V_m(0)=\frac4{a_{|m|}-1},\qquad
 V_{m,T}(0)=-\frac4{a_{|m|}+1}.
\]
In particular $V_{\pm1}=2$ identically, whereas
$0<V_m(0)\le1/4$ for $|m|\ge2$ and $V_{m,T}(0)<0$.
The section $T$ itself has $V=-2$.
\end{proposition}
\begin{proof}
At $w=0$ the curve is the nodal cubic
\[
                         y^2=x(x-2)^2.
\]
Its nonsingular affine locus is parametrized by
\[
 x=q^2,\qquad y=q(q^2-2).
\]
On this locus the chord law gives
\[
 q(R+S)=\frac{q(R)q(S)+2}{q(R)+q(S)}.
\]
Equivalently, over $\Q(\sqrt2)$ the coordinate
$(q+\sqrt2)/(q-\sqrt2)$ is multiplicative. At $P$ it equals
$3+2\sqrt2$. It follows that
\[
 q([m]P)=\frac{2b_m}{a_m-1},\qquad
 x([m]P)=\frac{2(a_m+1)}{a_m-1},\qquad
 y([m]P)=\frac{8b_m}{(a_m-1)^2}.
\]
All these points are nonsingular; hence the same group-law
computation gives the regular specialization of the rational
sections. The formula is unchanged in $x$ when $m$ is replaced by
$-m$. Translation by $T$ sends $x$ to $4(w+1)/x$, so its
specialization gives the second formula. Finally $a_1=3$,
$a_2=17$, and $a_m$ is strictly increasing.
\end{proof}

\section{The actual preparation}
For the conditional detector retain its actual tests $C_i$ and
\[
 F(t)=\frac{t(t+1)(2t+1)(3t+1)(4t+1)(5t+2)}{120},\qquad
 w(t)=\frac{F(t)}{4(t-1)^8},
\]
\[
 W=\prod_{i=1}^{6}\left(1+\frac{C_i^2}{F(t)}\right),
 \qquad V_{\mathrm{act}}=\Tr_{\Q(\sqrt2)/\Q}(W^{-1}).
\]
The retained real preparation places the points in fixed compact
neighborhoods on which all test functions are defined and bounded.
Both embeddings have $F(t)>0$. Since $C_1=1/x(H)$ is nonzero,
both real values of $W$ are strictly greater than one. Consequently
\[
 0<V_{\mathrm{act}}<2,\qquad
 w(t)\longrightarrow0,\qquad
 V_{\mathrm{act}}\longrightarrow2\quad(t\longrightarrow+\infty),
\]
uniformly on these neighborhoods.

\begin{corollary}\label{cor:tail}
Fix a finite set of integer multiples of $P$, allowing translation
by $T$. In the stated conditional preparation, for all sufficiently
large $t$ none of their trace functions equals $V_{\mathrm{act}}$.
The same holds for both arguments $t_j=a_j\lambda$, with the
actual fixed positive scales $a_j$, as $\lambda\longrightarrow+\infty$.
\end{corollary}
\begin{proof}
Every section in the finite set other than $\pm P$ has a limiting
trace different from two by Proposition~\ref{prop:limits}.
The exceptional sections $\pm P$ have trace exactly two and are
excluded by the strict inequality $V_{\mathrm{act}}<2$.
\end{proof}

\begin{proposition}\label{prop:contact}
Fix a finite set of multiples $[m]P$ with $|m|\ge2$ and of
translates $[m]P+T$. There is a finite prime set $S_0$, determined
by these indices, such that no corresponding trace equals an actual
detector trace at a fresh contact $p\notin S_0$ where
\[
            w\in p\Z_p,\qquad V_{\mathrm{act}}\equiv1
                    \ \text{or}\ 2\pmod p.
\]
For a section $[m]P$, it suffices to exclude the prime divisors of
$2(a_{|m|}-1)(a_{|m|}-3)(a_{|m|}-5)$.
For $[m]P+T$, it suffices to exclude the prime divisors of
$2(a_{|m|}+1)(a_{|m|}+3)(a_{|m|}+5)$; the section $T$ itself
requires only a fixed finite exclusion.
\end{proposition}
\begin{proof}
Outside two the special fiber at $w\equiv0$ is nodal, and $P,T$
are in its nonsingular locus. The same chord-law calculation
reduces the multiple to the fractions in
Proposition~\ref{prop:limits}, whenever the displayed denominators
are units. For the first fraction, its differences from one and
two have numerators $5-a_{|m|}$ and $2(3-a_{|m|})$.
For the translated fraction the corresponding numerators are
$-(a_{|m|}+5)$ and $-2(a_{|m|}+3)$. The exclusions make these
differences units. The original section $T$ has trace $-2$ and
is treated directly.
\end{proof}

For example, the duplication trace is
\[
                  V_2(w)=\frac{w^2-6w+1}{4}.
\]
It can have residue one or two at a contact only at the exceptional
odd primes three or seven. The original prime-selection theorem
permits arbitrary finite avoidance, so this section cannot match
the original prepared contact family.

\section{What a varying index would require}
At a split odd prime choose $\sqrt2\in\mathbb F_p$ and put
$r=3+2\sqrt2$. Whenever $r^m\ne1$, the same specialization gives
\[
 V_m(0)=\frac{8r^m}{(r^m-1)^2}.
\]
Thus trace residue two is equivalent to
\[
                    m\equiv\pm1\pmod{\operatorname{ord}(r)}.
\]
If also $\sqrt3\in\mathbb F_p$, trace residue one is equivalent to
\[
               r^m\in\{5+2\sqrt6,\,5-2\sqrt6\}.
\]
These are necessary residue conditions, not rational realization
theorems. They allow indices depending on the selected primes;
neither proposition excludes such indices or individual rational
points on the actual incidence cover. Introducing an unrestricted
index or another free elliptic point would require a separate
coordinate count. The present results preserve the established
counts and establish no smaller defining polynomial.
\end{document}
